\section{Introduction and position}
\label{sec:intro}

Dense retrieval finds paragraphs that \emph{look like} the question. A multi-hop question also
needs paragraphs that do not: the second document of a HotpotQA bridge question is tied to the
first by a fact, not by wording~\citep{yang2018hotpotqa}. The question that guided this work
was how little structure Dense retrieval needs to recover that evidence. The aim was never to
build a graph system. It was to find the cheapest structure that measurably helps, and to
measure its cost, so that richer structure has a baseline it must justify itself against.

\paragraph{Position.} The systems in this paper are \emph{untrained, call no language model at
query time, and are cheap}. We place them under that description and nowhere else. They are not
offered as competitors of trained multi-hop retrievers or of retrieval loops that call a language
model, and Section~\ref{sec:published} says what a reader may and may not conclude from the
published figures of those systems.

\paragraph{What we did.} A frozen first stage (a small Dense model, BM25, and the Entity Hop of
Section~\ref{sec:hop}) was measured on three benchmarks: HotpotQA~\citep{yang2018hotpotqa},
MuSiQue~\citep{trivedi2022musique} and MultiHop-RAG~\citep{tang2024multihoprag}. Every weight
was fitted once, on HotpotQA development questions, and then frozen. The last phase adds a
reordering step on top of the stored rankings: three zero-shot judges score every candidate the
four systems retrieved, nothing is fitted, and the paper reports whether the hop's candidates
still add evidence once a judge orders them.

\paragraph{Contributions.}
\begin{enumerate}
  \item A one-hop, query-blind Entity Hop whose indexing cost we measure on three corpora,
        with a zero-shot extractor that keeps most of the gain of a language-model extractor
        for a small fraction of its cost (Section~\ref{sec:hop}).
  \item Its measured gain over Dense and over Dense plus BM25 at full HotpotQA scale, its
        transfer to MuSiQue with nothing refitted, and its \emph{regression} on news, reported
        with the same test on every benchmark (Section~\ref{sec:scale}).
  \item The first measurement in this line of a zero-shot cross-encoder, the cheap standard
        control it had been missing, in two sizes, plus an open bi-encoder decision model,
        over all four stored systems and the union of their candidates, with throughput and
        dollar cost beside every figure: the strong cross-encoder lifts every system, the
        hop system with candidate relevance stays the best of the four under it on the two
        Wikipedia benchmarks, and the decision model lowers every system
        (Section~\ref{sec:judges}).
  \item The metrics the literature reports, computed beside this line's metric of record, and a
        table of figures published by other systems on the same questions, every row read in
        the paper it comes from (Section~\ref{sec:published}).
  \item Negative results, kept: concept spaces induced from pooled embeddings add nothing
        (Section~\ref{sec:concepts}); hard filters on the hop's seed entities lost on
        development data three times (Section~\ref{sec:hop}); the bi-encoder judge hurts on
        every set (Section~\ref{sec:judges}).
\end{enumerate}

\paragraph{Reading the numbers.} Every figure in the text, tables and plots is written by a
script from a committed artifact (Appendix~\ref{app:repro}). A figure is \emph{measured} when a
run produced it, \emph{derived} when it is arithmetic on measured values, and a
\emph{projection} when it is neither; projections are labelled where they appear.
